\documentclass[11pt,a4paper]{article}

% ---- Fonts and page ---------------------------------------------------------
\usepackage[T1]{fontenc}
\usepackage[utf8]{inputenc}
\usepackage{libertine}              % Linux Libertine text, Biolinum as the sans-serif font
\usepackage[libertine]{newtxmath}   % maths that matches Libertine
\usepackage[scaled=0.85]{beramono}  % monospaced font, scaled to match the x-height
\usepackage[a4paper,left=3.2cm,right=3.2cm,top=3cm,bottom=3.2cm]{geometry}
% Wide margins keep lines near 70 characters, which is what makes a page calm to read.
\usepackage{microtype}              % better spacing and hyphenation; costs nothing
\linespread{1.05}
\setlength{\parindent}{1.2em}
\setlength{\parskip}{0pt}

% ---- Maths, figures, tables ------------------------------------------------
\usepackage{amsmath}   % newtxmath already provides the amssymb symbols
\usepackage{graphicx}
\usepackage{booktabs}
\usepackage[font=small,labelfont=bf,labelsep=period,margin=0.6cm]{caption}
\usepackage{pgfplots}               % figures drawn from data, so the template needs no image files
\pgfplotsset{compat=1.16}

% ---- Colour, headings, page headers ----------------------------------------
\usepackage[dvipsnames]{xcolor}
\definecolor{accent}{HTML}{1F4E79}  % one accent colour, used for headings and links only
\usepackage{titlesec}
\titleformat{\section}{\sffamily\Large\bfseries\color{accent}}{\thesection}{0.7em}{}
\titleformat{\subsection}{\sffamily\large\bfseries}{\thesubsection}{0.7em}{}
\titlespacing*{\section}{0pt}{2.2ex plus .4ex}{1.2ex}
\titlespacing*{\subsection}{0pt}{1.8ex plus .3ex}{0.8ex}
\usepackage{fancyhdr}
\pagestyle{fancy}
\fancyhf{}
\renewcommand{\headrulewidth}{0pt}
\fancyhead[L]{\small\sffamily\color{gray} Preprint}
\fancyhead[R]{\small\sffamily\color{gray} A Short Running Title}
\fancyfoot[C]{\small\thepage}
\fancypagestyle{plain}{% first page: no running head, a status line instead
  \fancyhf{}\renewcommand{\headrulewidth}{0pt}
  \fancyfoot[L]{\footnotesize\sffamily\color{gray} Preprint, not peer reviewed}
  \fancyfoot[R]{\small\thepage}}

% ---- Authors and affiliations ---------------------------------------------
% authblk numbers the affiliations and lists them under the authors.
\usepackage[noblocks]{authblk}
\renewcommand{\Authfont}{\large}
\renewcommand{\Affilfont}{\small\itshape}
\renewcommand{\Authsep}{, }
\renewcommand{\Authand}{, }
\renewcommand{\Authands}{, and }
\setlength{\affilsep}{0.8em}

% ---- Citations and links ----------------------------------------------------
% Author-year citations: \citet{key} gives "Smith (2020)", \citep{key} gives "(Smith, 2020)".
\usepackage[round,authoryear]{natbib}
\bibliographystyle{plainnat}
\setlength{\bibsep}{3pt}
% hyperref goes last. Coloured links look the same on screen and on arXiv's PDF;
% they are not printed in colour by most printers, so nothing is lost on paper.
\usepackage[colorlinks=true,linkcolor=accent,citecolor=accent,urlcolor=accent,
            bookmarksnumbered]{hyperref}
\usepackage[capitalise]{cleveref}

\title{\sffamily\bfseries A Descriptive Title That Names the Method and the Result}
\author[1,*]{First Author}
\author[1,2]{Second Author}
\author[2]{Third Author}
\author[3]{Fourth Author}
\affil[1]{Department of Computer Science, Example University, City, Country}
\affil[2]{Example Research Institute, City, Country}
\affil[3]{Independent researcher}
\affil[*]{Corresponding author: \texttt{first.author@example.org}}
\date{\today}
% The PDF's title and author fields (what arXiv, search engines and reference managers
% read) are set by hand: hyperref's pdfusetitle cannot read authblk's author list.
\hypersetup{pdftitle={A Descriptive Title That Names the Method and the Result},
  pdfauthor={First Author, Second Author, Third Author, Fourth Author},
  pdfkeywords={keyword one, keyword two, keyword three}}

% The abstract in a narrower block with the keywords underneath.
\renewenvironment{abstract}{\begin{quote}\small\noindent\textbf{\sffamily Abstract.}\ }%
  {\end{quote}}
\newcommand{\keywords}[1]{\par\smallskip\noindent\textbf{\sffamily Keywords:} #1}

\begin{document}
\maketitle
\thispagestyle{plain}

\begin{abstract}
Summarise the paper in about 150 to 250 words: the problem, why it is hard, what you did,
and the one result a reader should remember. Avoid citations and equations here, because
repositories and search engines show the abstract on its own.
\keywords{keyword one, keyword two, keyword three}
\end{abstract}

\section{Introduction}
Open with the problem and why it matters. Then say what is missing in earlier work
\citep{he2016,vaswani2017} and what this paper adds. Author-year citations read well
in running text: \citet{kingma2015} is a good example of a method paper whose name carries
the idea. End the section with a short list of contributions and a roadmap
(\cref{sec:method} describes the method, \cref{sec:results} the results).

\section{Method}\label{sec:method}
Define the setting and notation before using it. Equations are numbered when you refer to them
later, so label only those you cite:
\begin{equation}
  \hat{\theta} = \arg\min_{\theta}\; \frac{1}{n}\sum_{i=1}^{n} \ell\bigl(f_{\theta}(x_i),\, y_i\bigr)
  + \lambda\,\lVert\theta\rVert_2^2 .
  \label{eq:objective}
\end{equation}
\Cref{eq:objective} is typeset with \texttt{amsmath} and cited with \texttt{\textbackslash cref},
which writes ``Equation~1'' for you.

\subsection{Implementation details}
Give enough detail to reproduce the work, and move the rest to Appendix~\ref{app:details}.
Mention software versions, random seeds and hardware.

\section{Results}\label{sec:results}
\Cref{fig:curve} shows a figure drawn directly in \LaTeX{} from data points;
\cref{tab:main} shows a table. Put the interpretation in the text, not only in the caption.

Report how much results vary, not just their mean: give the spread over runs or seeds
(standard deviation or a confidence interval) and say how many runs it is based on. If you
claim that one method is better, say which test supports it, or that the difference is within
the noise.

\begin{figure}[ht]
  \centering
  \begin{tikzpicture}
    \begin{axis}[width=0.86\linewidth, height=5.6cm, xlabel={Training steps (thousands)},
                 ylabel={Validation error (\%)}, xmin=0, xmax=50, ymin=0,
                 grid=major, grid style={gray!25}, legend pos=north east,
                 legend style={font=\small, draw=none, fill=none},
                 tick label style={font=\small}, label style={font=\small}]
      \addplot[thick, color=gray, dashed] coordinates {(0,38) (5,22) (10,17) (20,14.2) (30,13.1) (40,12.7) (50,12.4)};
      \addplot[thick, color=accent] coordinates {(0,36) (5,18) (10,12.5) (20,9.6) (30,8.6) (40,8.2) (50,8.1)};
      \legend{Baseline, Proposed}
    \end{axis}
  \end{tikzpicture}
  \caption{Validation error during training. Captions go below figures, and should be
  readable without the main text.}
  \label{fig:curve}
\end{figure}

\begin{table}[ht]
  \centering
  \caption{Test error (\%) of the baseline and the proposed method, mean over five runs.
  Captions go above tables.}
  \label{tab:main}
  \begin{tabular}{lcc}
    \toprule
    Method & Error (\%) & Time per step (ms) \\
    \midrule
    Baseline & 12.4 & 38 \\
    Proposed & \textbf{8.1} & 41 \\
    \bottomrule
  \end{tabular}
\end{table}
% Tables use only horizontal rules (booktabs): no vertical lines, no double rules.

\section{Discussion}
Interpret the results, state the limitations plainly, and say what a reader should not
conclude from them. A short conclusion can follow here or get its own section.

\section*{Acknowledgements}
Thank funders (with grant numbers), colleagues and reviewers. Delete the section if you have no one to thank.

\bibliography{refs}

\appendix
\section{Additional experimental details}\label{app:details}
Material that supports the paper but would interrupt it goes here: hyperparameters, proofs,
extra figures and tables. \Cref{tab:hyper} lists the settings used for \cref{sec:results}.

\begin{table}[ht]
  \centering
  \caption{Hyperparameters used in all experiments.}
  \label{tab:hyper}
  \begin{tabular}{ll}
    \toprule
    Setting & Value \\
    \midrule
    Optimiser & Adam \citep{kingma2015} \\
    Learning rate & $10^{-3}$ \\
    Batch size & 128 \\
    Weight decay $\lambda$ & $10^{-4}$ \\
    \bottomrule
  \end{tabular}
\end{table}

\section{Notes on submitting to arXiv}
arXiv builds your sources itself, and it does not run BibTeX. Upload the \texttt{.bbl} file that
Lemmata writes next to the PDF (\texttt{build/main.bbl}) together with \texttt{main.tex}, and any figure
files; keep the name of the \texttt{.bbl} the same as the main file. Delete this section before you submit.

\end{document}
